% Appendix A --- from operators to kernels, and the deferred proofs of section 4. Moved out of
% section 4 (Q-0291, the author's decision of 2026-10-04 to shorten sections 3 and 4 by moving
% material to appendices, as Q-0290 did for section 5; records/affine/
% move-sections-3-4-to-appendices*.py, three passes).
%   - lem:operators-to-kernels moved whole, with its `% shared with blueprint` marker, verbatim, with
%     the paragraph on its converses and the one on the operator form of thm:affine-rigidity; no
%     later argument of section 4 consumes it, and 4.1 keeps a paragraph saying what it gives.
%   - the proofs, unchanged but for their headings, in main-text order: prop:isotropic-corner, with
%     the two paragraphs naming what it uses (4.1); lem:affine-gauge, lem:parameter-space (4.2);
%     thm:cascade-infinitely-divisible (4.3); lem:charfn-one-on-subspace (with the sentences on a
%     negative-definite form of it), lem:affine-reduction, lem:single-ratio-cauchy (4.4). Section 4
%     keeps the proofs of thm:rank-deficient-step and thm:affine-rigidity.
%   - rem:no-covariance: the second paragraph of rem:kelker (4.5), verbatim (paper-only text).
% Q-0309 (fix group 2; records/affine/move-q0309-flow.py, move-q0309-no-covariance.py): this is
% Appendix A, the first of the proof appendices. It gains the proofs of lem:gw-kernels-from-origin
% and prop:gw-ray-families, whose statements moved to section 3.1, and loses rem:no-covariance,
% which became open problem 7 of section 5 (review flag 38).
% Q-0312 (review flag 20; records/affine/move-q0312-corner-parameter-space.py): prop:isotropic-corner
% with its two paragraphs, and the body of section 3.2 (lem:affine-gauge, lem:parameter-space and
% their prose), moved here verbatim, each before its proof; the main text keeps one paragraph for
% each. Edited after the move: the defences of [Iso]'s terminology (flags 10, 11) cut, [Iso] cited by
% statement name.

\section{Between operators and kernels, and the axiom scheme}\label{app:classical}

The two lemmas of the first subsection pass between the operator form of the scheme of
\S\ref{sec:affine-covariance} and its kernel form, in both directions. The other subsections prove the statements of
\S\S\ref{sec:affine-covariance}--\ref{sec:affine-single-step} whose proofs the main text defers, in
their order. Two of them also state what the main text summarizes in a paragraph: the isotropic
corner, and the two lemmas on the parameter space.

\subsection{Between operators and kernels}

The axioms of the scheme are stated on operators, in Definition~\ref{def:gw-cascade-family}(1), as the
line paper states its own (Definition~\ref{def:line-family}), and the paper works on kernels. The two
lemmas of this subsection show that nothing is lost. Lemma~\ref{lem:operators-to-kernels} passes from a
family of operators to its kernels, as Lemma~\ref{lem:line-representation} does on the line, and
Lemma~\ref{lem:kernels-to-operators} passes back.
% Round 3 (content review), B.R2 (VERIFY-round3-front-classical-paper.md; READING-wendel1952left.md,
% READING-larsen1971introduction.md; Wendel's p. 252 re-read at the page image for this edit).
Clause (1) of the first is, on the real space $L^1(\RR^d)$, Wendel's theorem on left centralizers
(\sscite{wendel1952left}, Thm.~1, p.~252; for operators commuting with the translations,
\sscite{larsen1971introduction}, Thm.~0.1.1, p.~2), with positivity and mass making the measure a
probability measure.

% shared with blueprint lem:operators-to-kernels
\begin{lemma}[from operators to kernels in dimension $d$]\label{lem:operators-to-kernels}
Fix $d \ge 1$. On the real space $L^1(\RR^d)$ write $T_af := f(\cdot - a)$ for $a \in \RR^d$,
$D_Af := |\det A|^{-1}f(A^{-1}\,\cdot)$ for $A \in \GL(d)$, and, for a probability measure $\mu$
on $\RR^d$, $(\mu * f)(x) := \int f(x - y)\,\mu(dy)$. Let $\varphi(x) := (2\pi)^{-d/2}e^{-|x|^2/2}$
be the standard Gaussian density.
\begin{enumerate}
\item \emph{(Representation.)} Let $\Phi : L^1(\RR^d) \to L^1(\RR^d)$ be linear and satisfy the
single-operator axioms of \restated{V}{Def.~3.1}{def:line-family} other than the reflection axiom (A3), read in $\RR^d$:
(A1) $\Phi$ is
bounded; (A2) $\Phi T_a = T_a\Phi$ for every $a \in \RR^d$; (A4) $\Phi f \ge 0$ whenever
$f \ge 0$; (A5) $\int\Phi f = \int f$ whenever $f \ge 0$. Then there is a unique probability
measure $\mu$ on $\RR^d$ with
\[
  \Phi f = \mu * f, \qquad f \in L^1(\RR^d) .
\]
Uniqueness holds in the strong form: if $\mu$ and $\nu$ are probability measures on $\RR^d$ with
$\mu * \varphi = \nu * \varphi$, then $\mu = \nu$. Conversely, for every probability measure $\mu$
on $\RR^d$ the operator $f \mapsto \mu * f$ satisfies (A1), (A2), (A4) and (A5).
\item \emph{(Transport.)} For every probability measure $\mu$ on $\RR^d$, every $A \in \GL(d)$ and
every $f \in L^1(\RR^d)$, $D_A(\mu * f) = (A_*\mu) * D_Af$. Consequently, for probability measures
$\mu$ and $\nu$, the operators satisfy $D_A\,(\mu * \cdot\,)\,D_A^{-1} = \nu * \cdot$ if and only
if $A_*\mu = \nu$.
\item \emph{(Continuity.)} If $\mu_n$ and $\mu$ are probability measures on $\RR^d$ with
$\mu_n * \varphi \to \mu * \varphi$ in $L^1(\RR^d)$, then $\widehat\mu_n(\omega) \to
\widehat\mu(\omega)$ for every $\omega \in \RR^d$.
\item \emph{(The $(G, W)$ families, and the affine family.)} Let $(G, W)$ be a wedge datum in
dimension $d$ and $(\Phi_{\Sigma,\Sigma'})_{(\Sigma,\Sigma') \in P_W}$ a $(G, W)$-cascade family of
operators (Definition~\ref{def:gw-cascade-family}(1)). Then the measures $\mu_{\Sigma,\Sigma'}$ that (1) gives,
$\Phi_{\Sigma,\Sigma'}f = \mu_{\Sigma,\Sigma'} * f$, form a $(G, W)$-cascade family of kernels
(Definition~\ref{def:gw-cascade-family}(2)). In particular, let $(\Phi_{\Sigma,\Sigma'})$ be linear
operators on $L^1(\RR^d)$, indexed by the pairs $\Sigma \le \Sigma'$ in $\PSD(d)$, with the axioms
of the case $(\GL(d), \text{Loewner})$: each $\Phi_{\Sigma,\Sigma'}$ satisfies (A1), (A2),
(A4) and (A5); (A6) $\Phi_{\Sigma,\Sigma} = \mathrm{Id}$ and
$\Phi_{\Sigma',\Sigma''}\Phi_{\Sigma,\Sigma'} = \Phi_{\Sigma,\Sigma''}$ whenever
$\Sigma \le \Sigma' \le \Sigma''$; (A7) for every $f \in L^1(\RR^d)$ the map
$(\Sigma,\Sigma') \mapsto \Phi_{\Sigma,\Sigma'}f$ is continuous from the set of pairs into
$L^1(\RR^d)$; $D_A\Phi_{\Sigma,\Sigma'}D_A^{-1} = \Phi_{A\Sigma A^{\tp},\,A\Sigma'A^{\tp}}$ for
every $A \in \GL(d)$ and every pair; and $\Phi_{0,I} \ne \mathrm{Id}$. Then the measures form an
affine cascade family in dimension $d$ (Definition~\ref{def:affine-cascade-family}).
\item \emph{(The isotropic family.)} Let $(\Phi_{s,t})_{0 \le s \le t}$ be linear operators on
$L^1(\RR^d)$ with these axioms: each $\Phi_{s,t}$ satisfies (A1),
(A2), (A4) and (A5); (A6) $\Phi_{t,t} = \mathrm{Id}$ and $\Phi_{s,t}\Phi_{r,s} = \Phi_{r,t}$ for
$r \le s \le t$; (A7) for every $f \in L^1(\RR^d)$ the map $(s,t) \mapsto \Phi_{s,t}f$ is
continuous from $\{0 \le s \le t\}$ into $L^1(\RR^d)$; $\Phi_{s,t}R_Q = R_Q\Phi_{s,t}$ for every
$Q \in \Orth(d)$, with $R_Qf := f(Q^{-1}\,\cdot)$; (A8) there are increasing bijections
$S_\lambda$ of $[0,\infty)$, $\lambda > 0$, with
$D_\lambda\Phi_{s,t} = \Phi_{S_\lambda s,\,S_\lambda t}D_\lambda$, where
$D_\lambda f := \lambda^{-d}f(\cdot/\lambda)$; and (ND) $\Phi_{s,t} \ne \mathrm{Id}$ whenever
$s < t$. Then the measures $\mu_{s,t}$ that (1) gives form an isotropic cascade family in
dimension $d$ (Definition~\ref{def:isotropic-kernel-family}).
\end{enumerate}
\end{lemma}

\begin{proof}
Two facts are used throughout. For a probability measure $\mu$ and $f \in L^1(\RR^d)$, $\mu * f$
is the $L^1(\RR^d)$-valued Bochner integral $\int T_yf\,\mu(dy)$: the map $y \mapsto T_yf$ is
continuous because translation acts continuously on $L^1(\RR^d)$, and it is bounded by
$\|f\|_1$. And with $\widehat h(\omega) := \int e^{i\langle\omega,x\rangle}h(x)\,dx$ for
$h \in L^1(\RR^d)$, Fubini gives the pairing
\[
  \widehat{\mu * f}(\omega) = \widehat\mu(\omega)\,\widehat f(\omega),
  \qquad |\widehat h(\omega)| \le \|h\|_1 ,
\]
and $\widehat\varphi(\omega) = e^{-|\omega|^2/2} > 0$.

\emph{(1), the interchange identity.} For $f, g \in L^1(\RR^d)$, $f * g = \int f(y)\,T_yg\,dy$ as
a Bochner integral, dominated by $|f(y)|\,\|g\|_1$. A bounded operator commutes with an
$L^1(\RR^d)$-valued integral, so (A1) and (A2) give $\Phi(f * g) = \int f(y)\,T_y\Phi g\,dy
= f * \Phi g$.

\emph{The approximants.} Let $B_\varepsilon := \{|x| < \varepsilon\}$, $|B_\varepsilon|$ its
volume, $\rho_\varepsilon := |B_\varepsilon|^{-1}\mathbf 1_{B_\varepsilon}$ and
$h_\varepsilon := \Phi\rho_\varepsilon$. By (A4) and (A5), $h_\varepsilon \ge 0$ with
$\int h_\varepsilon = 1$, so $\nu_\varepsilon := h_\varepsilon\,dx$ is a probability measure on
$\RR^d$. For every $f$, $\|f * \rho_\varepsilon - f\|_1 \le \sup_{|y| < \varepsilon}
\|T_yf - f\|_1 \to 0$ as $\varepsilon \downarrow 0$.

\emph{Tightness.} By the interchange identity $\rho_1 * h_\varepsilon = \Phi(\rho_1 *
\rho_\varepsilon)$, which tends to $\Phi\rho_1$ in $L^1(\RR^d)$ by (A1). Along
$\varepsilon_n := 1/(n+1)$ the functions $\rho_1 * h_{\varepsilon_n}$ therefore have uniformly
small tails: given $\eta > 0$ there is $R$ with $\int_{|x| > R}(\rho_1 * h_{\varepsilon_n}) < \eta$
for every $n$ (choose $N$ with $\|\rho_1 * h_{\varepsilon_n} - \Phi\rho_1\|_1 < \eta/2$ for
$n \ge N$, then $R$ for $\Phi\rho_1$ and the first $N$ terms). Convolution with a nonnegative
unit-mass density carried by the unit ball moves mass by at most $1$: for $|y| > R + 1$ the
translate $\rho_1(\cdot - y)$ is carried by $\{|x| > R\}$, so by Tonelli
\[
  \int_{|x| > R}(\rho_1 * h_\varepsilon)(x)\,dx
  = \int h_\varepsilon(y)\int_{|x| > R}\rho_1(x - y)\,dx\,dy
  \ge \nu_\varepsilon\bigl(\{|y| > R + 1\}\bigr) .
\]
So the sequence $(\nu_{\varepsilon_n})$ is tight, with the closed balls of radius $R + 1$ as
compacts.

\emph{The limit.} By Prokhorov's theorem a subsequence $\nu_{\varepsilon_{n_k}}$ converges weakly
to a probability measure $\mu$ on $\RR^d$.

\emph{Identification.} Let $\Psi \in L^\infty(\RR^d)$, acting on $L^1(\RR^d)$ by
$\Psi(h) := \int\Psi h$, and $f \in L^1(\RR^d)$. The map $y \mapsto \Psi(T_yf)$ is continuous and
bounded by $\|\Psi\|_\infty\|f\|_1$, hence a test function for weak convergence. Since
$f * h_\varepsilon = \int h_\varepsilon(y)\,T_yf\,dy$ and $\Psi$ commutes with the integral,
\[
  \Psi(f * h_{\varepsilon_{n_k}}) = \int\Psi(T_yf)\,\nu_{\varepsilon_{n_k}}(dy)
  \ \xrightarrow{\ k \to \infty\ }\ \int\Psi(T_yf)\,\mu(dy) = \Psi(\mu * f) .
\]
On the other hand $f * h_\varepsilon = \Phi(f * \rho_\varepsilon) \to \Phi f$ in $L^1(\RR^d)$, by
the interchange identity and (A1). Hence $\Psi(\Phi f - \mu * f) = 0$ for every $\Psi$; taking
$\Psi$ the sign of $\Phi f - \mu * f$ gives $\Phi f = \mu * f$.

\emph{Uniqueness.} If $\mu * \varphi = \nu * \varphi$, the pairing gives
$\widehat\mu(\omega)\,\widehat\varphi(\omega) = \widehat\nu(\omega)\,\widehat\varphi(\omega)$, so
$\widehat\mu = \widehat\nu$ because $\widehat\varphi > 0$, and $\mu = \nu$ because a finite measure
on $\RR^d$ is determined by its Fourier transform. In particular two measures representing the
same operator agree.

\emph{Converse.} For a probability measure $\mu$, $\|\mu * f\|_1 \le \|f\|_1$ by Minkowski's
integral inequality (A1); $T_a(\mu * f) = \mu * T_af$, both being
$x \mapsto \int f(x - a - y)\,\mu(dy)$ (A2); $\mu * f \ge 0$ when $f \ge 0$ (A4); and
$\int\mu * f = \int f$ for $f \ge 0$ by Tonelli (A5).

\emph{(2).} For $x \in \RR^d$, substituting $z = Ay$,
\[
  \bigl((A_*\mu) * D_Af\bigr)(x) = \int |\det A|^{-1}f\bigl(A^{-1}(x - z)\bigr)\,(A_*\mu)(dz)
  = |\det A|^{-1}\int f(A^{-1}x - y)\,\mu(dy) = D_A(\mu * f)(x) .
\]
$D_A$ is invertible on $L^1(\RR^d)$ with inverse $D_{A^{-1}}$, so applying the identity to
$D_A^{-1}g$ gives $D_A\,(\mu * D_A^{-1}g) = (A_*\mu) * g$ for every $g$: the operator
$D_A\,(\mu * \cdot\,)\,D_A^{-1}$ is convolution by $A_*\mu$. It is convolution by $\nu$ if and
only if $A_*\mu = \nu$, by the uniqueness clause of (1).

\emph{(3).} By the pairing, $|\widehat\mu_n(\omega) - \widehat\mu(\omega)|\,\widehat\varphi(\omega)
\le \|\mu_n * \varphi - \mu * \varphi\|_1 \to 0$; divide by $\widehat\varphi(\omega) > 0$. The test
density must have a zero-free transform, which is why the indicator of a ball would not do.

\emph{(4).} By (1), each $\Phi_{\Sigma,\Sigma'}$ is convolution by a unique probability measure
$\mu_{\Sigma,\Sigma'}$. The four axioms of the kernel form follow.
\begin{itemize}
\item (K1): $\Phi_{\Sigma,\Sigma} = \mathrm{Id}$ is convolution by $\delta_0$, so
$\mu_{\Sigma,\Sigma} = \delta_0$ by uniqueness. By Fubini
$\mu_{\Sigma',\Sigma''} * (\mu_{\Sigma,\Sigma'} * f) = (\mu_{\Sigma',\Sigma''} *
\mu_{\Sigma,\Sigma'}) * f$, so the cascade says that $\mu_{\Sigma,\Sigma''}$ and
$\mu_{\Sigma',\Sigma''} * \mu_{\Sigma,\Sigma'}$ represent the same operator, and they are equal by
uniqueness.
\item (K2): $P_W$ is a subset of a finite-dimensional space, so continuity on it is sequential. If
$(\Sigma_n,\Sigma'_n) \to (\Sigma,\Sigma')$ in $P_W$, continuity at $f = \varphi$ gives
$\mu_{\Sigma_n,\Sigma'_n} * \varphi \to \mu_{\Sigma,\Sigma'} * \varphi$ in $L^1(\RR^d)$, and (3)
gives $\widehat\mu_{\Sigma_n,\Sigma'_n}(\omega) \to \widehat\mu_{\Sigma,\Sigma'}(\omega)$.
\item (K3): for $A \in G$, $A\cdot\Sigma \le_W A\cdot\Sigma'$ because $W$ is $G$-invariant, so the
right-hand side of the covariance axiom is a member of the family, and (2) with
$\mu = \mu_{\Sigma,\Sigma'}$ and $\nu = \mu_{A\cdot\Sigma,A\cdot\Sigma'}$ turns the axiom into
$A_*\mu_{\Sigma,\Sigma'} = \mu_{A\cdot\Sigma,\,A\cdot\Sigma'}$.
\item (K4): if $\mu_{0,I} = \delta_0$ then $\Phi_{0,I}f = \delta_0 * f = f$ for every $f$, against
$\Phi_{0,I} \ne \mathrm{Id}$.
\end{itemize}
The pair $(\GL(d), \text{Loewner})$ is a wedge
datum: $\GL(d)$ is closed in itself and contains the dilations, and the Loewner order is
$\GL(d)$-invariant, since $A(\Sigma' - \Sigma)A^{\tp}$ is positive semidefinite with
$\Sigma' - \Sigma$, and contains the rays, since $c'\Sigma - c\Sigma = (c' - c)\Sigma$; its kernel
form is Definition~\ref{def:affine-cascade-family}.

\emph{(5).} The same argument on the parameter set $\{0 \le s \le t\}$ gives (I1) from (A6) and
(I2) from (A7). (I3): $|\det Q| = 1$, so $R_Q = D_Q$, and the commutation reads
$D_Q\Phi_{s,t}D_Q^{-1} = \Phi_{s,t}$; (2) with $\nu = \mu = \mu_{s,t}$ gives
$Q_*\mu_{s,t} = \mu_{s,t}$. (I4): $D_\lambda = D_{\lambda I}$, and (A8) reads
$D_{\lambda I}\Phi_{s,t}D_{\lambda I}^{-1} = \Phi_{S_\lambda s,\,S_\lambda t}$; (2) gives
$(\lambda I)_*\mu_{s,t} = \mu_{S_\lambda s,\,S_\lambda t}$. (ND): if $\mu_{s,t} = \delta_0$ with
$s < t$ then $\Phi_{s,t} = \mathrm{Id}$, against (ND).
\end{proof}

The converses of (4) and (5) hold as well, and the next lemma states them: convolution
with the kernels of a $(G, W)$-cascade family is a $(G, W)$-cascade family of operators, and
convolution with an isotropic cascade family satisfies the operator axioms of (5). The step
that is not formal is continuity in $L^1$, which (K2) gives only pointwise in the frequency.
Clause (1) of the lemma supplies it by smoothing with a Gaussian, without weak convergence.

% shared with blueprint lem:kernels-to-operators
\begin{lemma}[from kernels to operators in dimension $d$]\label{lem:kernels-to-operators}
Fix $d \ge 1$, and use the notation of Lemma~\ref{lem:operators-to-kernels}.
\begin{enumerate}
\item \emph{(Continuity.)} If $\mu_n$ and $\mu$ are probability measures on $\RR^d$ with
$\widehat\mu_n(\omega) \to \widehat\mu(\omega)$ for every $\omega \in \RR^d$, then
$\mu_n * f \to \mu * f$ in $L^1(\RR^d)$ for every $f \in L^1(\RR^d)$.
\item \emph{(The $(G, W)$ families.)} Let $(G, W)$ be a wedge datum in dimension $d$ and
$(\mu_{\Sigma,\Sigma'})_{(\Sigma,\Sigma') \in P_W}$ a $(G, W)$-cascade family of kernels
(Definition~\ref{def:gw-cascade-family}(2)). Then the operators $\Phi_{\Sigma,\Sigma'}f := \mu_{\Sigma,\Sigma'} * f$
on $L^1(\RR^d)$ form a $(G, W)$-cascade family of operators (Definition~\ref{def:gw-cascade-family}(1)), and the
measures that Lemma~\ref{lem:operators-to-kernels}(1) gives for them are the $\mu_{\Sigma,\Sigma'}$. So, for
every wedge datum, convolution is a one-to-one correspondence between the $(G, W)$-cascade families of
kernels and those of operators, and its inverse is the passage of
Lemma~\ref{lem:operators-to-kernels}(4).
\item \emph{(The isotropic family.)} Let $(\mu_{s,t})_{0 \le s \le t}$ be an isotropic cascade family
in dimension $d$ (Definition~\ref{def:isotropic-kernel-family}). Then the operators $\Phi_{s,t}f := \mu_{s,t} * f$
on $L^1(\RR^d)$ satisfy the axioms of Lemma~\ref{lem:operators-to-kernels}(5), with the bijections
$S_\lambda$ of (I4) in (A8), and the measures that Lemma~\ref{lem:operators-to-kernels}(1) gives for them are
the $\mu_{s,t}$.
\end{enumerate}
\end{lemma}

\begin{proof}
For $\sigma > 0$ write $\varphi_\sigma(x) := \sigma^{-d}\varphi(x/\sigma)$, the centred Gaussian density
with covariance $\sigma^2I$. For a probability measure $\nu$ and $f, g \in L^1(\RR^d)$ three facts are
used, each by Fubini's theorem: $\|\nu * f\|_1 \le \|f\|_1$; $\|g * f\|_1 \le \|g\|_1\|f\|_1$; and
$\nu * (g * f) = (\nu * g) * f$, where $\nu * g$ is the integrable function
$x \mapsto \int g(x - y)\,\nu(dy)$.

\emph{(1), the smoothed measures.} The Gaussian integral gives, for every $z \in \RR^d$,
\[
  \varphi_\sigma(z) = (2\pi)^{-d}\int e^{-i\langle\omega,z\rangle}e^{-\sigma^2|\omega|^2/2}\,d\omega .
\]
Put $z = x - y$ and integrate in $\nu(dy)$. The integrand is bounded by $e^{-\sigma^2|\omega|^2/2}$,
which is integrable for $d\omega \otimes \nu$, so Fubini's theorem gives, for every probability
measure $\nu$ and every $x \in \RR^d$,
\[
  (\nu * \varphi_\sigma)(x) = (2\pi)^{-d}\int e^{-i\langle\omega,x\rangle}\,\widehat\nu(\omega)\,
  e^{-\sigma^2|\omega|^2/2}\,d\omega .
\]
Put $h_n := \mu_n * \varphi_\sigma$ and $h := \mu * \varphi_\sigma$. Since $|\widehat\mu_n| \le 1$ and
$\widehat\mu_n \to \widehat\mu$ pointwise, dominated convergence gives $h_n(x) \to h(x)$ for every
$x$. The $h_n$ and $h$ are nonnegative with integral $1$ (Tonelli), so $\int(h_n - h) = 0$ and
\[
  \|h_n - h\|_1 = 2\int(h - h_n)^+ \xrightarrow{\ n \to \infty\ } 0 ,
\]
by dominated convergence again, since $0 \le (h - h_n)^+ \le h$ and $(h - h_n)^+ \to 0$ pointwise.

\emph{(1), the approximation.} By Minkowski's integral inequality and the substitution $y = \sigma z$,
\[
  \|\varphi_\sigma * f - f\|_1 \le \int\varphi_\sigma(y)\,\|T_yf - f\|_1\,dy
  = \int\varphi(z)\,\|T_{\sigma z}f - f\|_1\,dz ,
\]
which tends to $0$ as $\sigma \downarrow 0$ by dominated convergence: the integrand is bounded by
$2\|f\|_1\varphi(z)$ and tends to $0$ pointwise, translation being continuous on $L^1(\RR^d)$. Now
split $f = (f - \varphi_\sigma * f) + \varphi_\sigma * f$. By the three facts,
$\mu_n * (\varphi_\sigma * f) = h_n * f$ and $\mu * (\varphi_\sigma * f) = h * f$, so
\[
  \|\mu_n * f - \mu * f\|_1 \le 2\,\|f - \varphi_\sigma * f\|_1 + \|(h_n - h) * f\|_1
  \le 2\,\|f - \varphi_\sigma * f\|_1 + \|h_n - h\|_1\,\|f\|_1 .
\]
Given $\varepsilon > 0$, choose $\sigma$ so that the first term is below $\varepsilon$, and then $n_0$ so
that the second is below $\varepsilon$ for $n \ge n_0$.

\emph{(2).} Each $\Phi_{\Sigma,\Sigma'}$ is linear and satisfies (A1), (A2), (A4) and (A5) by the
converse in Lemma~\ref{lem:operators-to-kernels}(1). \emph{Cascade (A6).} $\Phi_{\Sigma,\Sigma} = \delta_0 * \cdot
= \mathrm{Id}$ by (K1), and for $\Sigma \le_W \Sigma' \le_W \Sigma''$, by the associativity of
convolution (Fubini) and (K1),
$\Phi_{\Sigma',\Sigma''}\Phi_{\Sigma,\Sigma'}f = \mu_{\Sigma',\Sigma''} * (\mu_{\Sigma,\Sigma'} * f)
= (\mu_{\Sigma',\Sigma''} * \mu_{\Sigma,\Sigma'}) * f = \Phi_{\Sigma,\Sigma''}f$.
\emph{Continuity (A7).} $P_W$ is a subset of the finite-dimensional space of pairs of symmetric matrices,
so continuity on it is sequential. If $(\Sigma_n,\Sigma'_n) \to (\Sigma,\Sigma')$ in $P_W$, (K2) gives
$\widehat\mu_{\Sigma_n,\Sigma'_n}(\omega) \to \widehat\mu_{\Sigma,\Sigma'}(\omega)$ for every $\omega$,
and (1) gives $\Phi_{\Sigma_n,\Sigma'_n}f \to \Phi_{\Sigma,\Sigma'}f$ in $L^1(\RR^d)$ for every $f$.
\emph{(ND).} If $\Phi_{0,I} = \mathrm{Id}$, then $\mu_{0,I}$ and $\delta_0$ represent the same
operator, so $\mu_{0,I} = \delta_0$ by the uniqueness clause of Lemma~\ref{lem:operators-to-kernels}(1),
against (K4). \emph{Covariance.} For $A \in G$ and $(\Sigma,\Sigma') \in P_W$, the pair
$(A\cdot\Sigma, A\cdot\Sigma')$ lies in $P_W$ because $W$ is $G$-invariant, and (K3) says
$A_*\mu_{\Sigma,\Sigma'} = \mu_{A\cdot\Sigma,\,A\cdot\Sigma'}$; by Lemma~\ref{lem:operators-to-kernels}(2)
this is $D_A\Phi_{\Sigma,\Sigma'}D_A^{-1} = \Phi_{A\cdot\Sigma,\,A\cdot\Sigma'}$. \emph{The measures.}
$\Phi_{\Sigma,\Sigma'}$ is convolution by $\mu_{\Sigma,\Sigma'}$, so by the uniqueness clause of
Lemma~\ref{lem:operators-to-kernels}(1) the measure that clause gives is $\mu_{\Sigma,\Sigma'}$. \emph{The
correspondence.} Kernels to operators to kernels returns the kernels, by the last step. Operators to
kernels to operators returns the operators, since the measures of
Lemma~\ref{lem:operators-to-kernels}(4) satisfy $\Phi_{\Sigma,\Sigma'}f = \mu_{\Sigma,\Sigma'} * f$.

\emph{(3).} The same argument on the parameter set $\{0 \le s \le t\}$, a subset of $\RR^2$, gives the
four single-operator axioms, (A6) from (I1), (A7) from (I2) and (1), and (ND) at every $s < t$ from
(ND). For $Q \in \Orth(d)$, $|\det Q| = 1$, so $D_Q = R_Q$, and (I3) with
Lemma~\ref{lem:operators-to-kernels}(2) gives $R_Q\Phi_{s,t}R_Q^{-1} = \Phi_{s,t}$, that is
$\Phi_{s,t}R_Q = R_Q\Phi_{s,t}$. For $\lambda > 0$, $D_\lambda = D_{\lambda I}$, and (I4) with
Lemma~\ref{lem:operators-to-kernels}(2) gives $D_\lambda\Phi_{s,t}D_\lambda^{-1} =
\Phi_{S_\lambda s,\,S_\lambda t}$, which is (A8). The measures are the $\mu_{s,t}$, as in (2).
\end{proof}

The sufficiency statements of the paper on $(G, W)$-cascade families and on isotropic families are
stated on kernels, and by Lemma~\ref{lem:kernels-to-operators} each holds on operators as well. The
rigidity on operators is Theorem~\ref{thm:affine-rigidity-operators}. The families of one linear flow
(Proposition~\ref{prop:one-parameter-families}) lie outside the scheme and are stated on kernels only.

\subsection{A family is its kernels from scale zero}\label{app:classical-kernels}

% Round 3 (report B, part 3 (vii); VERIFY-round3-front-classical-paper.md, B.x): the proof of record
% transcribed verbatim, so that clause (3) reads the G-invariance of W for (K3), as the blueprint does.
\begin{proof}[Proof of Lemma~\ref{lem:gw-kernels-from-origin}]
\emph{(1).} $P_W \subseteq P_{W'}$, and a triple $\Sigma \le_W \Sigma' \le_W \Sigma''$ is a triple for
$W'$, so (K1) holds on $P_W$; (K2) is the restriction of a continuous map to the subspace $P_W$; (K3)
holds for every $A \in G \subseteq G'$, and $(A\cdot\Sigma, A\cdot\Sigma') \in P_W$ by the
$G$-invariance of $W$; and $(0, I) \in P_W$, so (K4) is the same condition.

\emph{(2).} (a) is (K1) at the pair $(0,0)$ and (K4). (b) is (K3) at the pair $(0, \Sigma)$, as
$A\cdot 0 = 0$. (c): $0 \le_W \Sigma$ for every $\Sigma$, since $W$ contains the rays, and
$\Sigma \mapsto (0, \Sigma)$ is continuous from $\PSD(d)$ into $P_W$; compose with (K2). (d) is (K1) at
$0 \le_W \Sigma \le_W \Sigma'$. If $\mu_{\Sigma'} = \mu_\Sigma * \rho$, then
$\widehat\rho = \widehat\mu_{\Sigma'}/\widehat\mu_\Sigma$ wherever $\widehat\mu_\Sigma \ne 0$, which is
everywhere when $\widehat\mu_\Sigma$ has no zero; so $\rho$ is determined by its transform.

\emph{(3).} By the last sentence, $\mu_{\Sigma,\Sigma'}$ is well defined and its transform is the
quotient $q_{\Sigma,\Sigma'} := \widehat\mu_{\Sigma'}/\widehat\mu_\Sigma$. At $\Sigma = 0$ the quotient is
$\widehat\mu_{\Sigma'}$, by (a), so $\mu_{0,\Sigma'} = \mu_{\Sigma'}$. (K1): $q_{\Sigma,\Sigma} = 1$, and
$q_{\Sigma',\Sigma''}\,q_{\Sigma,\Sigma'} = q_{\Sigma,\Sigma''}$ whenever $\Sigma \le_W \Sigma' \le_W \Sigma''$;
uniqueness of the transform turns these into $\mu_{\Sigma,\Sigma} = \delta_0$ and the cascade identity.
(K2): by (c) the numerator and the denominator of $q_{\Sigma,\Sigma'}(\omega)$ are continuous on
$\PSD(d) \times \PSD(d)$ and the denominator has no zero, so the quotient is continuous there, and in
particular on $P_W$. (K3): for $A \in G$, $(A\cdot\Sigma, A\cdot\Sigma') \in P_W$ by $G$-invariance, and by
(b) the transform of $A_*\mu_{\Sigma,\Sigma'}$ at $\omega$ is
$q_{\Sigma,\Sigma'}(A^{\tp}\omega) = \widehat\mu_{A\cdot\Sigma'}(\omega)/\widehat\mu_{A\cdot\Sigma}(\omega)
= q_{A\cdot\Sigma,A\cdot\Sigma'}(\omega)$. (K4): $\mu_{0,I} = \mu_I \ne \delta_0$.
\end{proof}


\subsection{The isotropic corner}

The isotropic paper states the isotropic theory with its own definition,
Definition~\ref{def:isotropic-kernel-family}. That definition differs from the scheme's isotropic
corner, read on the multiples of $I$, in two ways. It leaves the relabelling $S_\lambda$ of the
scales free, where the scheme lets the dilation $\lambda I$ act on the multiples of $I$ by
$s \mapsto \lambda^2 s$. And it requires nondegeneracy at every pair, where the scheme requires it
at $(0, I)$ alone. The next proposition shows that, when $G$ contains $\Orth(d)$, neither difference
matters on the slice: the two theories
agree up to a change of the scale coordinate, which is unique up to a constant. The new coordinate
is $\varphi = \chi^2$, where $\chi$ is the coordinate of \prog{V}{Prop.~6.6} in which the dilations act
by multiplication, $\chi(S_\lambda t) = \lambda\,\chi(t)$. In \prog{V}{Rem.~6.7} $\varphi$ is the
\emph{variance gauge}: in it the Gaussian member is a one-parameter semigroup, and its scale is
proportional to the variance. The proposition calls the restriction of a family to the pairs
$(sI, tI)$ its \emph{slice}. Its proof, at the end of this subsection, reads an isotropic family
through its marginal family on a line (Lemma~\ref{lem:isotropic-marginal-family}) and builds the
families it needs from the ray order (Proposition~\ref{prop:gw-ray-families}).
% CHANGED Q-0281 wave 2: blind review W2C-2: the list of what the proposition uses follows the
% wave-2 proof (no longer prop:similarity-ray-families(3)).

% shared with blueprint prop:isotropic-corner+b1800e754ff1
\begin{proposition}[the slice of the isotropic corner is the isotropic theory]\label{prop:isotropic-corner}
Let $d \ge 1$, and let $\mathrm S(d) := \{cQ : c > 0,\ Q \in \Orth(d)\}$ be the similarities. Let
$(G, W)$ be a wedge datum in dimension $d$ (Definition~\ref{def:gw-cascade-family}) and
$(\nu_{\Sigma,\Sigma'})$ a $(G, W)$-cascade family of kernels, with kernels from scale zero
$\nu_\Sigma := \nu_{0,\Sigma}$. Since $W$ contains the rays, $(sI, tI) \in P_W$ whenever
$0 \le s \le t$; the family $(\nu_{sI,tI})_{0 \le s \le t}$ is the \emph{slice} of
$(\nu_{\Sigma,\Sigma'})$.
\begin{enumerate}
\item \emph{(The slice is an isotropic family.)} If $\Orth(d) \subseteq G$, the slice of every
$(G, W)$-cascade family of kernels is an isotropic cascade family in dimension $d$
(Definition~\ref{def:isotropic-kernel-family}) with the dilation action $S_\lambda s = \lambda^2 s$. In
particular $\nu_{sI,tI} \ne \delta_0$ whenever $s < t$.
\item \emph{(Every isotropic family is a slice, in the variance gauge.)} Let $(\mu_{s,t})$ be an
isotropic cascade family in dimension $d$ with dilation action $(S_\lambda)$. There is a gauge
$\chi$ with $\chi(S_\lambda t) = \lambda\,\chi(t)$ for every $\lambda > 0$ and $t \ge 0$; put
$\varphi := \chi^2$, a gauge. Then there is exactly one $(\GL(d), \text{rays})$-cascade family of
kernels $(\nu_{\Sigma,\Sigma'})$ with
\[
  \mu_{s,t} = \nu_{\varphi(s)I,\,\varphi(t)I}, \qquad 0 \le s \le t ,
\]
and its restriction to the datum $(\mathrm S(d), \text{rays})$ is an $(\mathrm S(d), \text{rays})$-cascade family of kernels
with the same property.
\item \emph{(The change of gauge is unique up to a constant.)} In the setting of (2), let $\psi$ be
a gauge, $(G, W)$ a wedge datum, and $(\nu'_{\Sigma,\Sigma'})$ a
$(G, W)$-cascade family of kernels with $\mu_{s,t} = \nu'_{\psi(s)I,\,\psi(t)I}$ for
$0 \le s \le t$. Then $\psi = c\,\varphi$ for some $c > 0$. Conversely, for every $c > 0$ the family
$(\nu_{c^{-1}\Sigma,\,c^{-1}\Sigma'})$ is a $(\GL(d), \text{rays})$-cascade family of kernels with
this property for $\psi = c\,\varphi$.
\item \emph{(What the slice determines.)} A $(\GL(d), \text{rays})$-cascade family of kernels is
determined by its slice. For $d \ge 2$ an $(\mathrm S(d), \text{rays})$-cascade family of kernels
is not: the slice of every such family is the slice of two such families whose kernels from scale zero differ at every $\Sigma \in \PD(d)$ that is not a multiple of $I$.
\end{enumerate}
\end{proposition}

Clause (1) says that, when $G$ contains $\Orth(d)$, nondegeneracy at the one pair $(0, I)$ gives
it at every pair of the slice. Off the slice it does not follow. For $d \ge 2$ the similarities
with the ray order admit a family whose kernel from scale zero at $\Sigma$ is
$((\det\Sigma)^{1/(2d)}I)_*m$, for a rotation-invariant
self-decomposable law $m \ne \delta_0$ (Proposition~\ref{prop:similarity-ray-families}(3)). Its
kernels at the singular
$\Sigma$ are all $\delta_0$, so its increments along the ray of such a $\Sigma$ are trivial, and
(ND) at every pair would exclude it. Clauses (2) and (3) say that the free gauge costs nothing:
the change of coordinate is the variance gauge $\varphi = \chi^2$, with $\chi$ taken from the
marginal family on a line (Lemma~\ref{lem:isotropic-marginal-family}). Clause (4) says what the corner adds to its slice:
nothing for $\GL(d)$, the shapes of the apertures for the similarities.

% CHANGED Q-0281 wave 2: blind review W2C-2: the lead and the next two paragraphs follow the
% wave-2 proof, which builds the shape-blind family through prop:gw-ray-families(2), proves S(d)
% closed in place, and reads lem:parameter-space(2), instead of prop:similarity-ray-families(3).
What the proof takes from the ray order is this. By
Proposition~\ref{prop:gw-ray-families}(3), a $(\GL(d), \text{rays})$-cascade family is determined by its
kernel at $I$, and for every rotation-invariant self-decomposable law $m \ne \delta_0$ there is one
with kernel $m$ at $I$, whose kernels from scale zero are $\nu_\Sigma = (\Sigma^{1/2})_*m$. By
Proposition~\ref{prop:gw-ray-families}(2), for a wedge datum with the ray order the kernels from scale zero of a family are self-decomposable, and every map $\Sigma \mapsto \nu_\Sigma$ into
self-decomposable laws with $\nu_0 = \delta_0$, $\nu_I \ne \delta_0$, covariance under the group and
transforms continuous in $\Sigma$ is the family of kernels from scale zero of exactly one family. For
the similarities $\Sigma \mapsto ((\det\Sigma)^{1/(2d)}I)_*m$ is such a map, and it gives the second
family of clause (4). By Lemma~\ref{lem:parameter-space}(2), a rotation-invariant law other than
$\delta_0$ has the orthogonal group as its stabilizer, which is what separates the two families off
the multiples of $I$.

The proof also uses three elementary facts. The transform of a self-decomposable law has no zero, and
a linear image $A_*\mu$ of a self-decomposable law is self-decomposable, with the residuals
$A_*\rho_c$ (Proposition~\ref{prop:gw-ray-families}(1)). When
$\widehat\nu_\Sigma$ has no zero, $\nu_{\Sigma'} = \nu_\Sigma * \rho$ determines
$\widehat\rho = \widehat\nu_{\Sigma'}/\widehat\nu_\Sigma$, so the increment $\nu_{\Sigma,\Sigma'}$
is the only such law $\rho$ (Lemma~\ref{lem:gw-kernels-from-origin}(2)). And a family of the general
linear group with the ray order is also one of the similarities, with the same kernels, since (K1),
(K2) and (K4) do not involve the group and (K3) for the similarities is part of (K3) for $\GL(d)$; the
similarities with the ray order are a wedge datum, as the proof shows.

\begin{proof}[Proof of Proposition~\ref{prop:isotropic-corner} (sketch)]
A fact used twice is the argument of \prog{V}{Lem.~6.3}: \emph{if a law $m$ on $\RR^d$ satisfies
$(aI)_*m = m$ for one $a > 0$ with $a \ne 1$, then $m = \delta_0$.} Indeed
$\widehat m(\omega) = \widehat m(a^n\omega)$ for every $n \in \ZZ$, and letting $a^n \to 0$ gives
$\widehat m(\omega) = 1$.

\emph{(1).} Since $W$ contains the rays and $Q \cdot (sI) = sI$ for $Q \in \Orth(d)$, (K1), (K2)
and (K3) of Definition~\ref{def:gw-cascade-family} give (I1), (I2) and (I3) of
Definition~\ref{def:isotropic-kernel-family}. For
$\lambda > 0$, $(\lambda I) \cdot (sI) = \lambda^2 s\,I$, so (K3) at $\lambda I$ is (I4) with
$S_\lambda s = \lambda^2 s$. For (ND), (K3) with $A = \sqrt t\,I$ gives
$\nu_{tI} = (\sqrt t\,I)_*\nu_I$. If $\nu_{sI,tI} = \delta_0$ with $s < t$, (K1) gives
$\nu_{tI} = \nu_{sI}$. For $s = 0$ this makes $\nu_I = \delta_0$. For $s > 0$ it makes $\nu_I$
invariant under $(\sqrt{s/t}\,I)_*$, and the fact gives $\nu_I = \delta_0$. Both contradict (K4).

\emph{(2).} The marginal family of Lemma~\ref{lem:isotropic-marginal-family}, the images
$\kappa_{s,t}$ of the $\mu_{s,t}$ on a line through the origin acting by $f \mapsto \kappa_{s,t} * f$,
satisfies (A1)--(A8) and (ND) with the same $(S_\lambda)$. \prog{V}{Prop.~6.6}, restated in the part used:
\emph{for such a family, $\ell(\lambda) := S_\lambda 1$ is a continuous strictly increasing bijection of
$(0,\infty)$ with $\ell(1) = 1$, and $\chi := \ell^{-1}$, with $\chi(0) := 0$, is an increasing bijection of
$[0,\infty)$ with $\chi(S_\lambda t) = \lambda\,\chi(t)$.} With $\theta := \varphi^{-1}$ this gives
$S_\lambda\theta(u) = \theta(\lambda^2 u)$. Put $m := \mu_{0,\theta(1)}$. By (I3), (I4) and (ND),
$m$ is rotation invariant, not $\delta_0$, and $\mu_{0,\theta(s)} = (\sqrt s\,I)_*m$ for $s \ge 0$.
By (I1) it is self-decomposable, \ref{eq:self-decomposable}, with the residuals
$\mu_{\theta(c^2),\theta(1)}$. Let $(\nu_{\Sigma,\Sigma'})$ be the $(\GL(d), \text{rays})$-cascade
family with kernel $m$ at $I$. Then $\nu_{uI} = \mu_{0,\theta(u)}$, and since the transforms of its
kernels from scale zero have no zero, (I1) and the uniqueness of increments
(Lemma~\ref{lem:gw-kernels-from-origin}(2)) give
$\nu_{uI,vI} = \mu_{\theta(u),\theta(v)}$ for $u \le v$. With $u = \varphi(s)$ and $v = \varphi(t)$
this is the claim. A second such family has the same kernel $m$ at $I$, hence is the same family.
For every $d \ge 1$ the similarities $\mathrm S(d)$ form a subgroup containing the dilations, closed
in $\GL(d)$, as an invertible limit $A$ of similarities has $AA^{\tp} = kI$ with $k > 0$; so
$(\mathrm S(d), \text{rays})$ is a wedge datum, the ray order being a partial order that contains the rays and is invariant
under $\GL(d)$.
% CHANGED Q-0281 wave 2: blind review W2C-3: the wedge datum moved here from (4), mirroring the blueprint;
% the two \ref's added in (1) and (2) name nodes the blueprint proof \uses (lem:gw-scheme-claims is
% not printed in this paper).
As $\mathrm S(d) \subseteq \GL(d)$, (K1), (K2) and (K4) do not involve the group, and (K3) for
$\mathrm S(d)$ is part of (K3) for $\GL(d)$, the family is also an
$(\mathrm S(d), \text{rays})$-cascade family, with the same kernels.

\emph{(3).} Put $m' := \nu'_I$, which is not $\delta_0$ by (K4). As in (1),
$\nu'_{uI} = (\sqrt u\,I)_*m'$. At the pairs $(0, t)$ the two families give
$(\sqrt{\psi(t)}\,I)_*m' = \mu_{0,t} = (\sqrt{\varphi(t)}\,I)_*m$, and at $t_0 := \theta(1)$ this
reads $m = (\sqrt{\psi(t_0)}\,I)_*m'$. So $m'$ is invariant under
$((\varphi(t)\psi(t_0)/\psi(t))^{1/2}I)_*$ for every $t > 0$, and the fact gives
$\psi = \psi(t_0)\,\varphi$. The converse is a check of (K1)--(K4) for
$\nu_{c^{-1}\Sigma,\,c^{-1}\Sigma'}$, the ray order being invariant under $\Sigma \mapsto c^{-1}\Sigma$.

\emph{(4).} Two $(\GL(d), \text{rays})$-cascade families with the same slice have the same kernel
at $I$, and are equal. Let $d \ge 2$, and let $(\nu_{\Sigma,\Sigma'})$ be an
$(\mathrm S(d), \text{rays})$-cascade family and $m := \nu_I$. It is rotation invariant,
self-decomposable and not $\delta_0$, and $\nu_{sI} = (\sqrt s\,I)_*m$. Proposition~\ref{prop:gw-ray-families}
gives two families with the kernels from scale zero $(\Sigma^{1/2})_*m$ and
$((\det\Sigma)^{1/(2d)}I)_*m$. They have the kernel $(\sqrt s\,I)_*m$ at $sI$ as well, and by the
uniqueness of increments the three slices agree. If the two kernels agree at $\Sigma \in \PD(d)$,
then $v^{-1}\Sigma^{1/2}$, with $v := (\det\Sigma)^{1/(2d)}$, fixes $m$, so it is orthogonal by
Lemma~\ref{lem:parameter-space}(2), and $\Sigma = v^2I$.
\end{proof}

\subsection{The parameter space}

The two lemmas of \S\ref{sec:parameter-space} say why the kernels of the scheme at the positive
definite matrices are indexed by $\PD(d)$, and what it costs. The first says that covariance fixes the
parametrization up to one constant. Their proofs follow them.

% shared with blueprint lem:affine-gauge
\begin{lemma}[the affine gauge is a constant]\label{lem:affine-gauge}
Let $h : \PD(d) \to \PD(d)$ satisfy $h(A\Sigma A^{\tp}) = A\,h(\Sigma)\,A^{\tp}$ for every
$A \in \GL(d)$ and $\Sigma \in \PD(d)$. Then there is a $c > 0$ with $h(\Sigma) = c\,\Sigma$ for
every $\Sigma$. On the isotropic slice $\{sI : s > 0\}$ the dilation $x \mapsto \lambda x$,
$\lambda > 0$, which is $A = \lambda I$, acts by $s \mapsto \lambda^2 s$.
\end{lemma}

The group acts transitively on $\PD(d)$, and the parametrization there is forced up to one
constant. On the multiples of $I$, Proposition~\ref{prop:isotropic-corner}(2) and (3) identify
the parameter with the variance gauge, up to that constant. In the Gaussian family
\ref{eq:affine-gaussian} the parameter is the covariance matrix up to the factor $2b$.

The second lemma says why $\PD(d)$ at all. A law that no proper subspace carries has a compact
stabilizer, and a rotation-invariant law other than $\delta_0$ has exactly $\Orth(d)$. An orbit of
$\GL(d)$ through a full law can therefore be indexed covariantly by $\PD(d) = \GL(d)/\Orth(d)$
exactly when it
has a rotation-invariant member.
% Round 3 (content review), B.R3 (READING-meerschaert2001limit.md, section 4, pp. 33-34 read at the
% page images for this edit).
Clause (1) is a classical consequence of the convergence of types (\sscite{meerschaert2001limit},
Thm.~2.3.10, p.~33, where the symmetries may carry a shift), and (3)(b) has the form of their
Thm.~2.3.11, p.~34, on the compact subgroups of $\GL(d)$.

% shared with blueprint lem:parameter-space+be463c730309
\begin{lemma}[the parameter space: compact stabilizers and the isotropic member]\label{lem:parameter-space}
Let $d \ge 1$. For a probability measure $\mu$ on $\RR^d$ write
$\operatorname{Stab}(\mu) := \{A \in \GL(d) : A_*\mu = \mu\}$ and
$\operatorname{Orb}(\mu) := \{A_*\mu : A \in \GL(d)\}$, and call $\mu$ \emph{full} if
$\mu(V) < 1$ for every linear subspace $V \ne \RR^d$.
\begin{enumerate}
\item \emph{(Compact stabilizers.)} $\operatorname{Stab}(\mu)$ is a closed subgroup of $\GL(d)$. If
$\mu$ is full, $\operatorname{Stab}(\mu)$ is compact. So a law whose stabilizer is not compact is
carried by a proper linear subspace; for $d \ge 2$ it need not be $\delta_0$: the stabilizer of
$\delta_{e_1}$ contains every $A$ with $Ae_1 = e_1$.
\item \emph{(Rotation-invariant laws.)} If $Q_*\mu = \mu$ for every $Q \in \Orth(d)$ and
$\mu \ne \delta_0$, then $\mu$ is full and $\operatorname{Stab}(\mu) = \Orth(d)$.
\item \emph{(The isotropic member is the axiom of the parameter space $\PD(d)$.)} Let $\mu_0$ be
full. The following are equivalent:
(a) $\operatorname{Orb}(\mu_0)$ has a rotation-invariant member;
(b) $\operatorname{Stab}(\mu_0) = C\,\Orth(d)\,C^{-1}$ for some $C \in \GL(d)$;
(c) there is a map $\Sigma \mapsto \mu_\Sigma$ from $\PD(d)$ into $\operatorname{Orb}(\mu_0)$ with
$A_*\mu_\Sigma = \mu_{A\Sigma A^{\tp}}$ for every $A \in \GL(d)$ and $\Sigma \in \PD(d)$.
When they hold, every map as in (c) is a bijection of $\PD(d)$ onto $\operatorname{Orb}(\mu_0)$, and
for any two such maps $\mu$ and $\mu'$ there is a $c > 0$ with $\mu'_\Sigma = \mu_{c\Sigma}$ for
every $\Sigma \in \PD(d)$.
\item \emph{(An orbit without an isotropic member.)} Let $p_1, \dots, p_d > 0$ be distinct with
$p_1 + \dots + p_d = 1$, and $\mu_0 := \sum_jp_j\,\delta_{e_j}$. Then $\mu_0$ is full,
$\operatorname{Stab}(\mu_0) = \{I\}$, the map $A \mapsto A_*\mu_0$ is a bijection of $\GL(d)$ onto
$\operatorname{Orb}(\mu_0)$, and no member of $\operatorname{Orb}(\mu_0)$ is rotation invariant.
\item \emph{(In the scheme.)} Let $(G, W)$ be a wedge datum with $\Orth(d) \subseteq G$ and
$(\mu_{\Sigma,\Sigma'})$ a $(G, W)$-cascade family of kernels (Definition~\ref{def:gw-cascade-family}(2)). Then
$\mu_I$ is rotation invariant and full, $\operatorname{Stab}(\mu_I) = \Orth(d)$, and
$\Sigma \mapsto \mu_\Sigma$ is injective on $\{AA^{\tp} : A \in G\}$, which is $\PD(d)$ when
$G = \GL(d)$. Without $\Orth(d) \subseteq G$, (K4) need not make $\mu_I$ full: for $d \ge 2$, $G$
the invertible diagonal matrices and $W$ the Loewner order, the centred Gaussian laws
$\mu_{\Sigma,\Sigma'}$ with $\widehat\mu_{\Sigma,\Sigma'}(\omega) =
\exp[-\tfrac12(\Sigma' - \Sigma)_{11}\,\omega_1^2]$ form a $(G, W)$-cascade family of kernels whose
$\mu_I$ is carried by the first axis and has a stabilizer that is not compact.
\end{enumerate}
\end{lemma}

For $G = \GL(d)$, clause (5) makes the kernels from scale zero at $\PD(d)$ correspond one to one to
$\GL(d)/\Orth(d) = \PD(d)$. For a smaller group containing $\Orth(d)$ it gives injectivity on
$\{AA^{\tp} : A \in G\}$ only. Without an isotropic member, a full law can have a trivial
stabilizer, and its orbit is then in bijection with $\GL(d)$ itself, of dimension $d^2$ against
$d(d+1)/2$ (clause (4)).

\begin{proof}[Proof of Lemma~\ref{lem:affine-gauge}]
For $Q \in \Orth(d)$, $QQ^{\tp} = I$ gives $Q\,h(I)\,Q^{\tp} = h(I)$: the symmetric matrix
$H := h(I)$ commutes with every orthogonal matrix. Taking $Q$ the reflection
$I - 2e_ie_i^{\tp}$, the entry $(QHQ^{\tp})_{ik}$ is $-H_{ik}$ for $k \ne i$, so $H_{ik} = 0$ and
$H$ is diagonal; taking $Q$ the permutation matrix exchanging $e_i$ and $e_k$ gives
$H_{ii} = H_{kk}$. So $H = cI$, and $c > 0$ because $H \in \PD(d)$. For $\Sigma \in \PD(d)$ put
$A = \Sigma^{1/2}$: then $h(\Sigma) = h(AIA^{\tp}) = A(cI)A^{\tp} = c\,\Sigma$. Finally
$(\lambda I)(sI)(\lambda I)^{\tp} = \lambda^2 s\,I$, which is how $\lambda I$ moves the parameter of
Definition~\ref{def:affine-cascade-family}.
\end{proof}

\begin{proof}[Proof of Lemma~\ref{lem:parameter-space} (sketch)]
\emph{(1).} By uniqueness of the transform, $A_*\mu = \mu$ exactly when
$\widehat\mu(A^{\tp}\omega) = \widehat\mu(\omega)$ for every $\omega$, so $\operatorname{Stab}(\mu)$ is
a subgroup, closed by the continuity of $\widehat\mu$. Let $\mu$ be full. The function
$f(B) := \int\min(1, |Bx|)\,\mu(dx)$ is continuous, by dominated convergence, and has a minimum $c$
on the compact sphere $\{\|B\| = 1\}$; $c > 0$, since $f(B) = 0$ would make $\mu$ carried by the
proper subspace $\ker B$. For $A$ in the stabilizer and $t := \|A\|$, since $A_*\mu = \mu$,
\[
  c \le f(A/t) = \int\min\bigl(1, |y|/t\bigr)\,\mu(dy) ,
\]
which tends to $0$ as $t \to \infty$. The stabilizer is therefore bounded, and as it contains the
inverses, $|\det A|$ is bounded away from $0$ on it. A limit of its elements is then invertible, and
in it, as it is closed in $\GL(d)$; so it is compact. For $\delta_{e_1}$ the matrices
$I + t\,e_1e_2^{\tp}$, $t \in \RR$, fix $e_1$ and are unbounded.

\emph{(2).} If a proper subspace carried $\mu$, so would a hyperplane $u^\perp$, hence by rotation
invariance every hyperplane, and so would the intersection $\{0\}$ of the coordinate hyperplanes.
So $\mu$ is full, and by (1) its stabilizer $K$ is compact and contains $\Orth(d)$. For $A \in K$
with polar decomposition $A = PQ$, the factor $P = AQ^{-1} \in \PD(d)$ lies in $K$, and so do its
powers $P^n$, $n \in \ZZ$. An eigenvalue $p$ of $P$ then has $p^n$ bounded, so $p = 1$ and $A = Q$.

\emph{(3).} For $B \in \GL(d)$, $B_*\mu$ is full with $\mu$, and
$\operatorname{Stab}(B_*\mu) = B\operatorname{Stab}(\mu)B^{-1}$. So (a) gives (b) by (2). From
(b), $\nu := (C^{-1})_*\mu_0$ is rotation invariant, and $\mu_\Sigma := (\Sigma^{1/2})_*\nu$ is as in
(c): for $\Sigma' := A\Sigma A^{\tp}$ the matrix $\Sigma'^{-1/2}A\Sigma^{1/2}$ is orthogonal. From
(c), $\mu_I$ is a rotation-invariant member. For a map as in (c), covariance with
$A = \Sigma^{1/2}$ gives $\mu_\Sigma = (\Sigma^{1/2})_*\mu_I$, which is onto the orbit and, by (2),
one to one. For a second map $\mu'$, $h := \mu^{-1} \circ \mu'$ satisfies
$h(A\Sigma A^{\tp}) = A\,h(\Sigma)A^{\tp}$, and Lemma~\ref{lem:affine-gauge} gives $h(\Sigma) = c\,\Sigma$.

\emph{(4).} A proper subspace misses some $e_j$, so it has $\mu_0$-measure at most $1 - p_j$. An
$A$ that fixes $\mu_0$ fixes each atom, the masses being distinct, so $A = I$. A member $A_*\mu_0$
gives positive mass only to the $d$ points $Ae_j$, and mass $p_1$ to $x := Ae_1 \ne 0$. If it were
rotation invariant, every point of norm $r := |x|$, being $Qx$ for an orthogonal $Q$, would carry
that mass and be one of those $d$ points; but the $2d$ points $\pm r\,e_j$ have norm $r$.

\emph{(5).} For $Q \in \Orth(d) \subseteq G$, $Q \cdot 0 = 0$ and $Q \cdot I = I$, so (K3) and (K4)
of Definition~\ref{def:gw-cascade-family} make $\mu_I$ rotation invariant and not $\delta_0$, and (2)
applies. (K3) gives $\mu_{AA^{\tp}} = A_*\mu_I$ for $A \in G$, and $A_*\mu_I = B_*\mu_I$ forces
$B^{-1}A \in \Orth(d)$, hence $AA^{\tp} = BB^{\tp}$. In the example the exponents are additive in
$\Sigma' - \Sigma$ and continuous, a diagonal $A$ multiplies them by $a_1^2$, and
$\widehat\mu_{0,I}(e_1) = e^{-1/2} \ne 1$. The stabilizer of its $\mu_I$ contains the unbounded
matrices $\diag(1, t, 1, \dots, 1)$, $t > 0$.
\end{proof}

\subsection{Infinitely divisible increments}

\begin{proof}[Proof of Theorem~\ref{thm:cascade-infinitely-divisible}]
\emph{Small increments.} Under the hypotheses of (1), for every $\varepsilon > 0$ and $\eta > 0$
there is $\delta > 0$ with $\mu_{a,b}(\{|x| > \varepsilon\}) < \eta$ whenever $0 \le a \le b \le 1$
and $b - a < \delta$. Otherwise there are pairs $0 \le a_n \le b_n \le 1$ with $b_n - a_n \to 0$ and
$\mu_{a_n,b_n}(\{|x| > \varepsilon\}) \ge \eta$. For each $\omega$ the function
$(a,b) \mapsto \widehat\mu_{a,b}(\omega)$ is uniformly continuous on the compact triangle
$\{0 \le a \le b \le 1\}$ and equal to $1$ on its diagonal, so $\widehat\mu_{a_n,b_n}(\omega) \to 1$.
By \Levy's continuity theorem on $\RR^d$, $\mu_{a_n,b_n} \to \delta_0$ weakly, and the portmanteau inequality for
the closed set $\{|x| \ge \varepsilon\}$ bounds the upper limit of its measures by $0$, a
contradiction.

% CHANGED Q-0281 wave 2: blind review W2C-5: the partition has n + 1 steps, as in the blueprint's
% proof and the ledger's A6 (r_n = n + 1); and W2C-1: the continuity theorem is named on R^d.
\emph{(1).} Fix $s \le t$, and for $n \ge 1$ put $t_{n,k} := s + k(t - s)/(n+1)$ and
$\mu_{n,k} := \mu_{t_{n,k-1},\,t_{n,k}}$, $1 \le k \le n+1$. The cascade identity and the commutativity
of convolution give, by induction on the number of steps,
\[
  \mu_{s,t} = \mu_{n,1} * \mu_{n,2} * \dots * \mu_{n,n+1} \qquad \text{for every } n .
\]
\emph{The rows are infinitesimal.} Given $\varepsilon > 0$ and $\eta > 0$, take $\delta$ from the
small increments: once $(t - s)/(n+1) < \delta$, every $\mu_{n,k}(\{|x| > \varepsilon\})$ is below
$\eta$. \emph{Khintchine.} The $\mu_{n,k}$ form a null array with $r_n = n + 1$ entries in row $n$,
whose row products all equal $\mu_{s,t}$, so C1 with $b_n = 0$ makes $\mu_{s,t}$ infinitely
divisible.

\emph{(2).} For a $W$-chain $\zeta$ from $\Sigma$ to $\Sigma'$, the laws
$\nu_{u,v} := \mu_{\zeta(u),\zeta(v)}$, $0 \le u \le v \le 1$, satisfy the hypotheses of (1) by
(K1) and (K2) of Definition~\ref{def:gw-cascade-family}, since $(u,v) \mapsto (\zeta(u),\zeta(v))$
is continuous. And $\nu_{0,1} = \mu_{\Sigma,\Sigma'}$. The kernel $\mu_\Sigma$ is the increment of
the pair $0 \le_W \Sigma$.

\emph{(3).} (a) $\zeta(t) - \zeta(s) = (t - s)(\Sigma' - \Sigma)$ is positive semidefinite. (b)
Along each chain a later value is a multiple of an earlier one by a factor at least $1$. (c) By the
spectral theorem $M^s$ is positive definite, continuous in $s$, and $M^aM^b = M^{a+b}$. For
$s \le t$ put $B := \Sigma^{1/2}M^{s/2}$, so that $\zeta(s) = BB^{\tp}$ and
$\zeta(t) = BM^{t-s}B^{\tp}$. Then $Q := \zeta(s)^{-1/2}B$ is orthogonal, and
$\zeta(s)^{-1/2}\zeta(t)\zeta(s)^{-1/2} = QM^{t-s}Q^{\tp}$ lies in $\mathcal K$ by the two
closure properties. The three sets named are closed under orthogonal conjugation, and
$M \mapsto M^\tau$ raises every eigenvalue, the least one and the determinant to the power
$\tau \ge 0$, which preserves each defining inequality. For $M = e^X$, $\lambda_{\min}(M) =
e^{\lambda_{\min}(X)}$ and $\det M = e^{\tr X}$, which gives the description of $\mathcal K_\alpha$.
The Loewner and the ray order are partial orders directly. That the relation of $\mathcal K_\alpha$
% Round 3 (content review), B.S7: the order clause of 4.7(2) spends no cited fact.
is one is the first sentence of Proposition~\ref{prop:stable-step-cone}(2), which rests on no cited fact.

\emph{(4).} For $s \le t$ the laws $\nu_{u,v} := \mu_{s + u(t-s),\,s + v(t-s)}$ satisfy the hypotheses
of (1) by (I1) and (I2) of Definition~\ref{def:isotropic-kernel-family}, and
$\nu_{0,1} = \mu_{s,t}$.
\end{proof}

\subsection{The single step}

\begin{proof}[Proof of Lemma~\ref{lem:charfn-one-on-subspace}]
Let $U := \{\eta \in H : |\eta| < r\}$, with $r > 0$ so small that $\widehat\rho = 1$ on $U$. For
$\eta \in U$, $0 = \mathrm{Re}\bigl(1 - \widehat\rho(\eta)\bigr)
= \int\bigl(1 - \cos\langle\eta,x\rangle\bigr)\rho(dx)$ with a nonnegative integrand, so
$\langle\eta,x\rangle \in 2\pi\ZZ$ for every $x$ outside a $\rho$-null set $N_\eta$. Let
$D \subseteq U$ be countable and dense in $U$, and $N := \bigcup_{\eta\in D}N_\eta$, a $\rho$-null
set. For $x \notin N$ the linear map $\ell_x : H \to \RR$, $\eta \mapsto \langle\eta,x\rangle$, is
continuous and maps $D$ into the closed set $2\pi\ZZ$, hence maps $U$, which lies in the closure
of $D$, into $2\pi\ZZ$. As $U$ is connected, $\ell_x(U)$ is a connected subset of $2\pi\ZZ$, a
single point, and it contains $\ell_x(0) = 0$; so $\ell_x = 0$ on $U$, hence on $H$, which is the
union of the dilates of $U$. That is $x \in H^\perp$. Thus $\rho(\RR^d \setminus
H^\perp) \le \rho(N) = 0$, and for $\eta \in H$,
$\widehat\rho(\omega+\eta) = \int_{H^\perp}e^{i\langle\omega,x\rangle}e^{i\langle\eta,x\rangle}\rho(dx)
= \widehat\rho(\omega)$. For $H = \RR^d$ the complement is $\{0\}$, so $\rho = \delta_0$.
\end{proof}

\begin{proof}[Proof of Lemma~\ref{lem:affine-reduction}]
\emph{The kernel from scale zero at $I$ is radial and real.} For $Q \in \Orth(d)$, (K3) of
Definition~\ref{def:affine-cascade-family} with $\Sigma = 0$,
$\Sigma' = I$ gives $Q_*\mu_I = \mu_{QQ^{\tp}} = \mu_I$, so
$\widehat\mu_I(Q^{\tp}\omega) = \widehat\mu_I(\omega)$ and $\widehat\mu_I$ is constant on spheres.
Put $\vartheta(u) := \widehat\mu_I(\sqrt u\,e_1)$ for $u \ge 0$; then
$\widehat\mu_I(\omega) = \vartheta(|\omega|^2)$, $\vartheta$ is continuous because a characteristic
function is, and $\vartheta(0) = 1$. With $Q = -I$,
$\widehat\mu_I(\omega) = \widehat\mu_I(-\omega) = \overline{\widehat\mu_I(\omega)}$, so $\vartheta$ is
real, and $|\vartheta| \le 1$.

\emph{Every kernel from scale zero.} For $\Sigma \in \PD(d)$, (K3) with $A = \Sigma^{1/2}$ gives
$\mu_\Sigma = A_*\mu_I$, so $\widehat\mu_\Sigma(\omega) = \widehat\mu_I(\Sigma^{1/2}\omega)
= \vartheta(\langle\omega,\Sigma\omega\rangle)$. For $\Sigma \in \PSD(d)$ singular,
$\Sigma_n := \Sigma + n^{-1}I \in \PD(d)$ tends to $\Sigma$, and (K2) along the pairs
$(0,\Sigma_n) \to (0,\Sigma)$ with the continuity of $\vartheta$ gives
$\widehat\mu_\Sigma(\omega) = \lim_n\vartheta(\langle\omega,\Sigma_n\omega\rangle)
= \vartheta(\langle\omega,\Sigma\omega\rangle)$.

\emph{No zeros.} Fix $\omega$ and put $f(s,t) := \widehat\mu_{sI,tI}(\omega)$ for $0 \le s \le t$.
By (K2) $f$ is continuous; by (K1), and because the transform of a convolution is the product of
the transforms, $f(t,t) = 1$ and $f(r,t) = f(s,t)\,f(r,s)$ for $r \le s \le t$. Suppose
$f(0,t) = 0$ for some $t$, and let $t^*$ be the infimum of such $t$. As $f(0,\cdot)$ is continuous
with $f(0,0) = 1$, $t^* > 0$ and $f(0,t^*) = 0$. By continuity at $(t^*,t^*)$, where $f = 1$, there
is $s \in (0,t^*)$ with $f(s,t^*) \ne 0$, and $f(0,s) \ne 0$ because $s < t^*$; then
$f(0,t^*) = f(s,t^*)\,f(0,s) \ne 0$, a contradiction. Hence
$\vartheta(t|\omega|^2) = \widehat\mu_{tI}(\omega) \ne 0$ for all $t \ge 0$ and $\omega$, so $\vartheta$
has no zero on $[0,\infty)$. Being real and continuous with $\vartheta(0) = 1$, it is positive, and
$g := -\log\vartheta$ is continuous, nonnegative because $\vartheta \le 1$, and $g(0) = 0$. This is
\ref{eq:affine-reduction}.

\emph{Increments.} For $\Sigma \le \Sigma'$, (K1) gives $\mu_{\Sigma'} = \mu_{\Sigma,\Sigma'} *
\mu_\Sigma$, so $\widehat\mu_{\Sigma'} = \widehat\mu_{\Sigma,\Sigma'}\,\widehat\mu_\Sigma$, and
$\widehat\mu_\Sigma > 0$ gives \ref{eq:affine-increment}.

\emph{Nondegeneracy.} If $g \equiv 0$ then $\widehat\mu_I \equiv 1$ and
Lemma~\ref{lem:charfn-one-on-subspace} with $H = \RR^d$ gives $\mu_I = \delta_0$, against (K4).
\end{proof}

\begin{proof}[Proof of Lemma~\ref{lem:single-ratio-cauchy}]
\emph{A constant difference on a half-line.} Fix $\delta > 0$ and put $u_\delta := \delta/(a-1)$,
so that $a u_\delta = u_\delta + \delta$. For $x \ge u_\delta$, \ref{eq:single-ratio} at
$u = u_\delta$ and $w = x - u_\delta \ge 0$ reads
$g(x+\delta) - g(x) = g(u_\delta + \delta) - g(u_\delta) =: \kappa(\delta)$, independent of $x$ on
$[u_\delta,\infty)$.

\emph{A constant difference everywhere.} Fix $\delta > 0$ and $y > 0$. For every integer
$n \ge \delta/((a-1)y)$ the points $y + j\delta/n$, $0 \le j < n$, are at least
$y \ge u_{\delta/n}$, so the previous step with $\delta/n$ in place of $\delta$ gives
\[
  g(y+\delta) - g(y) = \sum_{j=0}^{n-1}\Bigl[g\bigl(y + \tfrac{(j+1)\delta}{n}\bigr)
     - g\bigl(y + \tfrac{j\delta}{n}\bigr)\Bigr] = n\,\kappa(\delta/n) .
\]
Given $y_1, y_2 > 0$, an $n$ admissible for both shows
$g(y_1+\delta) - g(y_1) = g(y_2+\delta) - g(y_2)$. So $K(\delta) := g(y+\delta) - g(y)$ does not
depend on $y > 0$, and $K(2\delta) = [g(y+2\delta) - g(y+\delta)] + [g(y+\delta) - g(y)]
= 2K(\delta)$.

\emph{Additivity, from continuity at $0$ alone.} Put $h(\delta) := g(\delta) - K(\delta)$ for
$\delta > 0$. With $y = \delta$ the previous step gives $g(2\delta) = g(\delta) + K(\delta)$, so
$h(2\delta) = g(\delta) + K(\delta) - 2K(\delta) = h(\delta)$. By induction, for every $n \ge 0$,
\[
  h(\delta) = h(\delta 2^{-n}) , \qquad K(\delta) = 2^nK(\delta 2^{-n}) , \qquad\text{hence}\qquad
  h(\delta) = g(\delta 2^{-n}) - 2^{-n}K(\delta) .
\]
Letting $n \to \infty$, continuity at $0$ and $g(0) = 0$ give $h(\delta) = 0$, that is
$g(\delta) = K(\delta)$. Hence $g(y + \delta) = g(y) + g(\delta)$ for all $y, \delta > 0$.

\emph{Cauchy's equation.} By induction $g(n\delta) = n\,g(\delta)$ for integers $n \ge 1$, so
$q\,g(p/q) = g(p) = p\,g(1)$ and $g(r) = r\,g(1)$ for every positive rational $r$. As $g \ge 0$,
$g(y+\delta) \ge g(y)$: $g$ is nondecreasing on $(0,\infty)$. For $x > 0$ and rationals
$0 < r < x < r'$ this gives $r\,g(1) \le g(x) \le r'\,g(1)$, and letting $r, r' \to x$ gives
$g(x) = x\,g(1)$. With $g(0) = 0$ that is the claim.
\end{proof}
